\chapter{Ethical Issues}
\label{ch:ethics}

\todo[color=blue!80!white, textcolor=white]{Rewrite ethics chapter}

This chapter considers the wider ethical, legal, and professional implications of the work. The primary ethical concerns posed by this project lie in the fact that the same feedback-guided attack strategies proposed to stress-test and improve system robustness could, in practice, be repurposed by malicious actors to evade detection in currently deployed systems.

\section{The Dual-Use Dilemma and Offensive Security}
The central contribution of this work is a closed-loop attack agent designed to degrade the performance of adaptive learning systems while evading state-of-the-art drift detectors \cite{Korycki_2022, Kuppa_2022}. By definition, this research falls under the umbrella of \textit{offensive security} or \textit{red teaming}.

This comes with an inherent risk that identifying the precise boundaries of current defences could provide a blueprint for attackers. However, it is widely accepted in the security community that "security through obscurity" (relying on the ignorance of attackers) is a fragile defence strategy. As noted in~\cref{ch:introduction}, adaptive systems are increasingly deployed in critical domains such as network intrusion detection and fraud prevention. If these systems harbour latent vulnerabilities to feedback-driven attacks, it is ethically preferable for these weaknesses to be identified and documented by this work rather than be discovered by potential real-world adversaries.

To mitigate the risks associated with this dual-use technology, this project adheres to the principle of defensive disclosure. The attack algorithms are also implemented solely within a sandboxed simulation environment using public datasets, ensuring no live systems are targeted.

\section{Legal and Professional Compliance}
From a legal perspective, this project is conducted in strict accordance with the \textit{Computer Misuse Act 1990 (UK)}. Section 3 of the Act prohibits unauthorised acts with the intent to impair the operation of a computer. 
\begin{itemize}
    \item \textbf{Authorisation :} All experiments are conducted on local, isolated hardware or authorised university compute clusters.
    \item \textbf{Target Scope:} The victim systems are local instantiations of open-source algorithms, not external services. At no point does this project interact with, probe, or attempt to subvert third-party servers or commercial APIs.
\end{itemize}

\section{Data Governance and Privacy}
Finally, the project relies exclusively on publicly available benchmark datasets. No personally identifiable information (PII) is collected, processed, or generated. The risk of privacy violation is therefore negligible. The focus remains strictly on the algorithmic dynamics of drift detection rather than the content of the data itself.